% --------------------------
% CONFIGURATION
% --------------------------
\documentclass[11pt,a4paper]{moderncv}
\moderncvstyle{classic}
\moderncvcolor{black}
\definecolor{firstnamecolor}{rgb}{0,0,0}

% Basic packages
\usepackage[utf8]{inputenc}
\usepackage[T1]{fontenc}
\usepackage[scale=0.8, top=1.5cm, bottom=1.5cm]{geometry}
\usepackage[dvipsnames]{xcolor}
\definecolor{PaleBlue}{RGB}{76, 151, 215}
\usepackage[colorlinks=true, urlcolor=PaleBlue, linkcolor=PaleBlue]{hyperref}
\usepackage{microtype}

% Configure font settings
\renewcommand{\rmdefault}{ppl}
\renewcommand{\sfdefault}{phv}
\renewcommand*{\namefont}{\fontsize{26}{18}\mdseries}

% Configure microtype
\microtypesetup{expansion=false}

% Load sensitive information
\input{../../secrets/personal_info_dk.tex}

% --------------------------
\begin{document}
\makecvtitle

Dear Professor Lioma,

\vspace*{2em}
I am applying for the PhD fellowship in Resource Efficiency for Generative AI at the University of Copenhagen. After some years in industry, I want to return to research and build a career researching machine learning. What I have enjoyed most in my work is investigating machine learning methods, examining their limitations, and discussing difficult questions with colleagues. I am excited to return to that work through research that can make capable AI less resource-intensive and more accessible.

\vspace{4pt}
\hspace*{2em}
My background in Engineering Physics, specialising in Statistical Learning and Data Analytics, provides a strong mathematical foundation for this work. At RaySearch Laboratories, I researched automated radiotherapy planning using deep learning, Gaussian processes, and numerical optimisation. This gave me experience investigating uncertainty and evaluating methodological choices. At Signum Life Science, I now lead the development of a forecasting platform supporting thousands of independently trained models and work on embedding-based text retrieval. I would bring experience designing reproducible experiments, investigating failures, and implementing the software needed to test research ideas.

\vspace{4pt}
\hspace*{2em}
For this opportunity, I am particularly interested in how generative AI systems can estimate whether further computation is likely to improve an answer. How can a system assess the difficulty of a problem without first solving it? This reminds me of the famous halting problem. My practical interest is in inexpensive, probabilistic estimates of the benefit of additional reasoning, and understanding when those estimates are reliable enough to guide decisions.

\vspace{4pt}
\hspace*{2em}
Building on existing work in adaptive stopping, I would like to investigate how stopping policies behave when tasks differ from those used for calibration. A starting point would be to study whether signals from intermediate outputs can reveal when a policy's estimates have become unreliable, and whether adjusting the stopping rule can reduce premature stops. Using tasks with verifiable answers, I would compare these policies with existing stopping methods and fixed computation budgets. Evaluation would cover answer quality, latency, and measured energy consumption, including the cost of monitoring and calibration.


\vspace*{2em}
The Inference \& Retrieval Lab's combination of expertise in machine learning, information retrieval, and efficient AI is especially appealing to me. I would value the opportunity to develop this direction through regular discussion with the supervisory team and fellow researchers, contributing my mathematical and engineering experience while growing as an independent researcher.


\vspace{2em}
Kind regards,

Ivar Cashin Eriksson

\end{document}
